\section{Scope, Method, and Qualifications}

\subsection{Reader, question, and document identity}

This is a discursive philosophical- and dogmatic-theological companion essay in \emph{The Creature Before God}. It is not a fourth movement added to the trilogy's progression from presence through due return to union. It addresses a distinct objection at the threshold between natural theology and received revelation: whether a universal God and universally significant truth can act and become known through determinate persons, events, words, and sacramental signs.

The intended reader is an intellectually serious Christian, or a reader considering the Christian claim, who can follow a metaphysical distinction but is not presumed to know scholastic terminology. The governing claim is that the recoil from historical and sacramental particularity often rests upon a confusion between an abstract universal and a universal cause. God is neither a general concept nor one member of a class. In the Christian account, the universal cause knows and sustains singular creatures, elects for mission, assumes a particular human nature in the Person of the eternal Son, and communicates the one Paschal mystery through determinate signs. The argument establishes the coherence and theological necessity of this pattern; it does not by itself prove that a reported historical event occurred.

The article is independently intelligible. Within the larger series, it develops the passage from the universal source and present dependence treated by \emph{Ontological Vertigo} to the revealed, ecclesial, and sacramental forms presupposed by \emph{The Due Return}.

\subsection{Terms and source functions}

In the article, \emph{particular} means concrete, singular, determinate, or historically located: this person, people, act, word, or sign. It never means that God is one specimen of divinity. An \emph{abstract universal} is a concept considered apart from individuating features; a \emph{universal cause} is a cause whose power extends to all the effects within its order. \emph{Election} means a gratuitous divine choice that gives a vocation and mission; it is not treated as proof of the recipient's antecedent superiority. \emph{Universal salvific scope} means that Christ's one mediation and the offer of salvation are ordered to all; it does not mean that all are in fact saved.

\begin{threestable}{Source class}{Function in the article}{Governing boundary}
Sacred Scripture & Supplies the pattern of Abrahamic promise, Israel's priestly and prophetic mission, the Incarnation, Christ's one mediation, material signs, the visible neighbor, and Christ's concluding question. & Brief quoted phrases remain in their canonical contexts; no verse is made to yield a later metaphysical distinction by itself.\\
Conciliar and catechetical teaching & Governs revelation through deeds and words, the historical Incarnation, Israel and the Church, respect for other religions, providence, universal salvific offer, and sacramental efficacy. & Distinct documents are read together; uniqueness does not become contempt, and universality does not become a claim that all are saved.\\
Authoritative doctrinal declaration & \emph{Dominus Iesus} supplies the bounded claim that Christ's finite human acts have the divine Person of the Word as their subject and that his mediation is unique and universal. & Its use is limited to nos.~6 and 13--15 and paired with \emph{Dei Verbum}, \emph{Fides et ratio}, \emph{Gaudium et spes}, and \emph{Nostra aetate}; its rhetoric is not made the article's comparative method.\\
Thomistic theology & Distinguishes God from a member of a genus or an abstract universal and accounts for divine knowledge of singulars. & The metaphysical distinction supports coherence; it neither demonstrates the Incarnation nor binds every permissible Catholic school to all Thomistic formulations.\\
Project synthesis & Names the ``flight into abstraction,'' orders election, Incarnation, sacrament, and neighbor under one argumentative problem, and supplies the literary contrasts. & These formulations are editorial synthesis, not formal ecclesial classifications or attributed definitions.\\
\end{threestable}

\subsection{Included and excluded scope}

Included are the difference between an abstract universal and universal cause; divine knowledge and providence of singulars; the outward purpose of Abraham's and Israel's election; the continuing Christian bond with the Jewish people; the Incarnation as a historical and personal event; the relation between Christ's particular human acts and universal salvific mediation; sacramental signs as embodied communication of grace; and the visible neighbor as the test of universal love.

Excluded are a complete proof of God, a historical defense of the Resurrection, a treatise on revelation or inspiration, a comprehensive theology of Israel and the Church, a comparative account of Buddhist or other non-Christian traditions, a full theology of religions, a treatise on sacramental validity and fruitfulness, and adjudication of individual salvation. The article does not classify its target as canonical or theological heresy. It does not defend philosophical nominalism, political particularism, nationalism, religious tribalism, or the claim that every alleged divine intervention is credible.

\subsection{Material qualifications}

\begin{enumerate}[leftmargin=1.45em,itemsep=0.25em]
\item God is not ``particular'' as a member of a genus. Particularity in this article concerns divine effects, address, the assumed human nature, historical acts, and sacramental means.
\item A universal concept and a universal cause are not competing theories of the same thing. The distinction prevents conceptual generality from being mistaken for ontological transcendence.
\item Apophatic theology and analogical predication remain necessary. God's self-revelation does not make the divine essence comprehensible or contained by created language.
\item Concrete providence neither eliminates secondary causality nor identifies every evil, illness, or suffering as a direct punishment or private divine message. Permission, natural processes, human agency, and moral responsibility remain real.
\item Election is gratuitous and vocational. The biblical movement from Abraham and Israel toward all nations does not make historical Israel dispensable, justify contempt, or settle every contemporary theological question about covenant and mission.
\item The article confesses Christian fulfillment in Christ while preserving Paul's prohibition against Gentile boasting and the Church's teaching that God's gifts and call to Israel are irrevocable.
\item Christ's human words and deeds are finite without being a merely partial divine appearance because their personal subject is the eternal Son. The divine mystery remains transcendent and inexhaustible.
\item Christ's unique and universal mediation does not imply that all persons are saved, that explicit historical knowledge is the only possible mode of receiving grace, or that truth and holiness are absent from other religions.
\item The Paschal mystery occurred once and is not historically repeated. Its sacramental presence follows the Church's account of Christ's abiding act, not a reversal of time.
\item Sacraments are neither magical techniques nor optional illustrations. Christ is their principal actor; the Church is bound to the received signs; proper disposition remains relevant to their fruit; God is not reduced to the visible instrument. In the Eucharist, the article confesses transubstantiation and Christ's true, real, and substantial presence without attempting a complete Eucharistic treatise.
\item Sacramental logic does not make every material event a sacrament in the strict sense or every encounter an immediate revelation.
\item Love of the visible neighbor tests claimed love of God without identifying creature and Creator by nature or reducing salvation to philanthropy.
\item Coherence is not historical proof. The article removes an antecedent objection to divine action in history; it does not exempt Christian claims from historical inquiry or the obedience of faith.
\item The final question is scriptural address, not a claim that the reader's response has been inferred, emotionally produced, or supplied by the article.
\end{enumerate}

\subsection{Translations, rights, currentness, and review}

Scriptural language is limited to short traditional or working-English phrases checked at the identified loci. Official conciliar, catechetical, and doctrinal English witnesses are cited at exact numbered sections. Thomas is cited by work, question, and article; the governing Latin corpus was checked through Corpus Thomisticum, with New Advent used as a public working-English aid. Scripture, creeds, official texts, source wording, and third-party translations retain their own rights status and are not relicensed by the project's CC BY 4.0 grant. The surrounding argument and editorial synthesis are project-created.

The identified online witnesses were checked through 23 July 2026. No mutable canon law, civil law, liturgical edition, territorial discipline, or jurisdiction-specific conclusion governs the article. Internal claim-to-source and rights review has been completed. Production and artifact evidence for the current source is recorded in the repository research scope. The installed snapshot remains on hold pending exact-snapshot authorization. Independent philosophical, biblical, Jewish--Christian, dogmatic, sacramental, interreligious, literary, and ecclesiastical review remains outstanding; no release clearance, imprimatur, nihil obstat, or other ecclesiastical approval is claimed.
